\chapter{Introduction}
\label{chap:introduction}

Marine vessels contain numerous brackets, foundations, lifting fittings, and control-surface structures. These components must transmit demanding loads without yielding or excessive deformation while contributing as little unnecessary mass as practicable. Excess mass increases material consumption and vessel displacement, whereas excessive material removal reduces stiffness and fatigue resistance. Marine structural design is therefore a constrained engineering problem.

\section{Problem Statement}
Conventional component design uses hand calculations, conservative initial sizing, CAD modelling, and repeated checks. Only a limited number of dimensional combinations can normally be examined manually. A systematic method was required to explore plate dimensions, stiffener sizes, and fillet radii while enforcing strength and deflection limits.

\section{Objectives}
The main objective was to develop and assess a parameterized engineering method for lower-mass marine structural components. The specific objectives were to:
\begin{enumerate}
\item define materials, loads, boundary conditions, variables, and acceptance criteria;
\item formulate analytical models for mass, stress, factor of safety, and deflection;
\item explore the bracket design space by Latin-hypercube sampling and bounded Nelder--Mead optimization;
\item generate parameterized geometry and validation candidates in Fusion 360;
\item extend the method to a lifting padeye and fixed-envelope stabilizer fin; and
\item identify the detailed FEA and physical validation required before practical adoption.
\end{enumerate}

\section{Scope and Significance}
Completed work covered linear-elastic analytical screening, continuous dimensional variables, CAD generation, and validation preparation. Fatigue, buckling, weld-detail stress, corrosion, nonlinear contact, vibration, impact, CFD, fabrication, and proof testing were outside the completed scope. The engineering value lies in traceable comparison of alternatives, not in replacing detailed analysis or professional judgement.
